\documentclass[journal=jacsat,manuscript=suppinfo]{achemso}

\usepackage[T1]{fontenc}
\usepackage[utf8]{inputenc}
\usepackage{amsmath}
\usepackage{amssymb}
\usepackage{booktabs}
\usepackage{multirow}
\usepackage{longtable}
\usepackage{array}
\usepackage{graphicx}
\usepackage{xcolor}
\usepackage{hyperref}
\usepackage{siunitx}
\usepackage{listings}
\usepackage{caption}

\lstset{
  basicstyle=\footnotesize\ttfamily,
  breaklines=true,
  frame=single,
  language=Python,
}

%% Author information (must match main.tex)
\author{Ku Kang}
\affiliation[CBDRI]{CBRN Defense Research Institute, Ministry of National
Defense, Seoul 06796, Republic of Korea}

\author{Jeongyun Kim}
\affiliation[SNUCBE]{Department of Chemical and Biological Engineering,
  Seoul National University, Seoul 08826, Republic of Korea}

\author{Jin Yoo}
\affiliation[CBDRI]{CBRN Defense Research Institute, Ministry of National
Defense, Seoul 06796, Republic of Korea}

\author{Dongyoul Lee}
\affiliation{Department of Physics and Chemistry, Korea Military
Academy, Seoul 01805, Republic of Korea}

\email{dlee32@kma.ac.kr}

\title{
Physics-Informed Machine Learning for Thermodynamic Stability Constants of
  Theranostic Radiometal--Chelator Complexes: Multi-Temperature Prediction
  via Physics-Constrained Van't~Hoff Neural Networks}

\begin{document}

\tableofcontents
\newpage

%%============================================================
%% TABLE S1: Descriptor Definitions
%%============================================================
\section*{Table S1. Physics-Informed Descriptor Definitions and Sources}
\label{tab:s1}

\begin{longtable}{@{}p{3.5cm}p{5.5cm}p{4cm}@{}}
  \toprule
  Descriptor & Definition & Source \\
  \midrule
  \endhead
  \multicolumn{3}{l}{\textit{Metal-centered descriptors}} \\
  Ionic radius & Shannon effective ionic radius (\AA) at coordination number 6 & Shannon (1976)\cite{Shannon1976} \\
  Charge & Formal oxidation state of metal ion & IUPAC \\
  HSAB hardness & Absolute chemical hardness $\eta$ (eV) & Pearson (1963)\cite{Pearson1963} \\
  Coord.\ number & Preferred coordination number in aqueous complexes & Literature \\
  \midrule
  \multicolumn{3}{l}{\textit{Chelator-centered descriptors}} \\
  Donor O & Number of carboxylate/hydroxamate oxygen donors & Ligand structure \\
  Donor N & Number of amine/imine nitrogen donors & Ligand structure \\
  Ring size & Macrocyclic cavity ring size (0 for linear) & Ligand structure \\
  Is macrocyclic & Binary indicator (1 = macrocyclic, 0 = acyclic) & Ligand structure \\
  \midrule
  \multicolumn{3}{l}{\textit{Metal--chelator descriptors}} \\
  Size mismatch & $\Delta r = r_\mathrm{ion} - r_\mathrm{cavity}$ (\AA) & Computed \\
  Cavity volume & Estimated cavity volume (\AA$^3$) & Computed from ring size \\
  log$K$ (EDTA ref) & Conditional stability constant of metal--EDTA at $I = 0.1$~M & NIST database \\
  \midrule
  \multicolumn{3}{l}{\textit{DFT descriptor (A2/A3 models only)}} \\
  $E_\mathrm{int}$ (DFT) & Interaction energy B3LYP-D3(BJ)/def2-TZVP//CPCM(water) (kJ~mol$^{-1}$) & Computed (22 pairs) \\
  \bottomrule
\end{longtable}

\newpage

%%============================================================
%% TABLE S2: Full Dataset
%%============================================================
\section*{Table S2. Full 140-Entry Dataset}
\label{tab:s2}

Conditional stability constants $\log K_1$ at $T = 25\,^\circ$C, $I = 0.1$~M.
Source codes: Exp = experimentally measured (potentiometry/ITC/spectrophotometry)
with primary literature citations; Est = estimated via linear free-energy
relationship (LFER) anchored to the NIST Critical Stability Constants
Database\cite{MartellSmith2004} following the methodology of
Hancock \& Martell.\cite{Hancock1989}

\begin{longtable}{@{}llcp{1.5cm}c@{}}
  \toprule
  Metal & Chelator & $\log K_1$ & Source & Ref. \\
  \midrule
  \endhead
  % Sc (12 entries)
  Sc & DOTA     & 30.79 & Exp & \cite{Pniok2014} \\
  Sc & NOTA     & 27.0  & Exp & \cite{Pniok2014} \\
  Sc & DTPA     & 24.3  & Est & \cite{MartellSmith2004,Hancock1989} \\
  Sc & EDTA     & 23.4  & Exp & \cite{Pniok2014} \\
  Sc & macropa  & 22.1  & Est & \cite{MartellSmith2004,Hancock1989} \\
  Sc & CB-TE2A  & 25.1  & Exp & \cite{Pniok2014} \\
  Sc & AAZTA    & 24.1  & Est & \cite{MartellSmith2004,Hancock1989} \\
  Sc & HOPO     & 25.3  & Est & \cite{MartellSmith2004,Hancock1989} \\
  Sc & Deferox  & 23.8  & Est & \cite{MartellSmith2004,Hancock1989} \\
  Sc & PCTA     & 30.5  & Exp & \cite{Pniok2014} \\
  Sc & TRAP     & 27.2  & Est & \cite{MartellSmith2004,Hancock1989} \\
  Sc & H4octapa & 23.5  & Est & \cite{MartellSmith2004,Hancock1989} \\
  % Bi (11 entries)
  Bi & DOTA     & 30.3  & Exp & \cite{Milne2015} \\
  Bi & NOTA     & 22.0  & Est & \cite{MartellSmith2004,Hancock1989} \\
  Bi & DTPA     & 35.0  & Exp & \cite{Milne2015} \\
  Bi & EDTA     & 21.0  & Exp & \cite{MartellSmith2004} \\
  Bi & macropa  & 24.5  & Exp & \cite{Morgan2023} \\
  Bi & AAZTA    & 30.0  & Est & \cite{MartellSmith2004,Hancock1989} \\
  Bi & HOPO     & 22.2  & Est & \cite{MartellSmith2004,Hancock1989} \\
  Bi & H4octapa & 34.9  & Exp & \cite{Morgan2023} \\
  Bi & CHX-DTPA & 27.8  & Est & \cite{MartellSmith2004,Hancock1989} \\
  Bi & PCTA     & 29.2  & Est & \cite{MartellSmith2004,Hancock1989} \\
  Bi & H4neunpa & 21.8  & Est & \cite{MartellSmith2004,Hancock1989} \\
  % Ga (19 entries)
  Ga & NOTA     & 31.0  & Exp & \cite{Notni2010} \\
  Ga & DOTA     & 21.3  & Exp & \cite{Notni2010} \\
  Ga & DTPA     & 25.5  & Exp & \cite{MartellSmith2004} \\
  Ga & EDTA     & 8.5   & Exp & \cite{MartellSmith2004} \\
  Ga & macropa  & 17.0  & Est & \cite{MartellSmith2004,Hancock1989} \\
  Ga & AAZTA    & 26.7  & Exp & \cite{Notni2010} \\
  Ga & TRAP     & 24.5  & Exp & \cite{Simecek2012} \\
  Ga & Deferox  & 20.3  & Exp & \cite{MartellSmith2004} \\
  Ga & HOPO     & 27.3  & Est & \cite{MartellSmith2004,Hancock1989} \\
  Ga & CB-TE2A  & 19.5  & Est & \cite{MartellSmith2004,Hancock1989} \\
  Ga & PCTA     & 30.8  & Exp & \cite{Notni2010} \\
  Ga & H4octapa & 21.7  & Est & \cite{MartellSmith2004,Hancock1989} \\
  Ga & H4neunpa & 21.9  & Est & \cite{MartellSmith2004,Hancock1989} \\
  Ga & TACN     & 10.3  & Exp & \cite{MartellSmith2004} \\
  Ga & TETA     & 20.1  & Est & \cite{MartellSmith2004,Hancock1989} \\
  Ga & CHX-DTPA & 24.2  & Est & \cite{MartellSmith2004,Hancock1989} \\
  Ga & NODA-GA  & 18.5  & Exp & \cite{Notni2010} \\
  Ga & DATA5m   & 21.0  & Exp & \cite{Simecek2012} \\
  Ga & THP      & 20.5  & Exp & \cite{Notni2010} \\
  % In (15 entries)
  In & DOTA     & 23.9  & Exp & \cite{MartellSmith2004} \\
  In & NOTA     & 26.2  & Exp & \cite{MartellSmith2004} \\
  In & DTPA     & 29.0  & Exp & \cite{MartellSmith2004} \\
  In & EDTA     & 15.7  & Exp & \cite{MartellSmith2004} \\
  In & macropa  & 19.2  & Est & \cite{MartellSmith2004,Hancock1989} \\
  In & AAZTA    & 18.0  & Exp & \cite{MartellSmith2004} \\
  In & H4octapa & 22.6  & Exp & \cite{Morgan2023} \\
  In & CB-TE2A  & 19.0  & Est & \cite{MartellSmith2004,Hancock1989} \\
  In & CHX-DTPA & 20.8  & Est & \cite{MartellSmith2004,Hancock1989} \\
  In & HOPO     & 26.0  & Est & \cite{MartellSmith2004,Hancock1989} \\
  In & PCTA     & 25.0  & Est & \cite{MartellSmith2004,Hancock1989} \\
  In & TRAP     & 25.2  & Est & \cite{MartellSmith2004,Hancock1989} \\
  In & H4neunpa & 15.5  & Est & \cite{MartellSmith2004,Hancock1989} \\
  In & NODA-GA  & 29.2  & Exp & \cite{Morgan2023} \\
  In & THP      & 18.2  & Est & \cite{MartellSmith2004,Hancock1989} \\
  % Zr (2 entries)
  Zr & DOTA     & 43.1  & Exp & \cite{Blei2023} \\
  Zr & DFO*     & 43.0  & Exp & \cite{Blei2023} \\
  % Lu (16 entries)
  Lu & DOTA     & 25.0  & Exp & \cite{deLeonRodriguez2008} \\
  Lu & EDTA     & 14.0  & Exp & \cite{MartellSmith2004} \\
  Lu & DTPA     & 22.4  & Exp & \cite{MartellSmith2004} \\
  Lu & NOTA     & 12.0  & Exp & \cite{MartellSmith2004} \\
  Lu & macropa  & 13.8  & Est & \cite{MartellSmith2004,Hancock1989} \\
  Lu & AAZTA    & 22.8  & Exp & \cite{Baranyai2007} \\
  Lu & HOPO     & 18.5  & Est & \cite{MartellSmith2004,Hancock1989} \\
  Lu & CB-TE2A  & 22.0  & Exp & \cite{MartellSmith2004} \\
  Lu & H4octapa & 19.8  & Exp & \cite{Morgan2023} \\
  Lu & PCTA     & 25.1  & Exp & \cite{MartellSmith2004} \\
  Lu & TRAP     & 22.4  & Est & \cite{MartellSmith2004,Hancock1989} \\
  Lu & CHX-DTPA & 18.3  & Est & \cite{MartellSmith2004,Hancock1989} \\
  Lu & H4neunpa & 19.6  & Est & \cite{MartellSmith2004,Hancock1989} \\
  Lu & NODA-GA  & 22.2  & Est & \cite{MartellSmith2004,Hancock1989} \\
  Lu & TACN     & 22.6  & Est & \cite{MartellSmith2004,Hancock1989} \\
  Lu & TETA     & 22.1  & Est & \cite{MartellSmith2004,Hancock1989} \\
  % Y (17 entries)
  Y  & DOTA     & 24.9  & Exp & \cite{MartellSmith2004} \\
  Y  & EDTA     & 13.0  & Exp & \cite{MartellSmith2004} \\
  Y  & DTPA     & 22.1  & Exp & \cite{MartellSmith2004} \\
  Y  & NOTA     & 11.0  & Exp & \cite{MartellSmith2004} \\
  Y  & macropa  & 12.7  & Est & \cite{MartellSmith2004,Hancock1989} \\
  Y  & AAZTA    & 22.1  & Exp & \cite{Baranyai2007} \\
  Y  & HOPO     & 17.7  & Est & \cite{MartellSmith2004,Hancock1989} \\
  Y  & CB-TE2A  & 21.3  & Exp & \cite{MartellSmith2004} \\
  Y  & H4octapa & 18.1  & Exp & \cite{Morgan2023} \\
  Y  & PCTA     & 24.7  & Exp & \cite{MartellSmith2004} \\
  Y  & TRAP     & 22.3  & Est & \cite{MartellSmith2004,Hancock1989} \\
  Y  & CHX-DTPA & 17.5  & Est & \cite{MartellSmith2004,Hancock1989} \\
  Y  & H4neunpa & 13.2  & Est & \cite{MartellSmith2004,Hancock1989} \\
  Y  & NODA-GA  & 22.0  & Est & \cite{MartellSmith2004,Hancock1989} \\
  Y  & TACN     & 11.2  & Est & \cite{MartellSmith2004,Hancock1989} \\
  Y  & TETA     & 21.5  & Est & \cite{MartellSmith2004,Hancock1989} \\
  Y  & DATA5m   & 14.5  & Est & \cite{MartellSmith2004,Hancock1989} \\
  % Ac (8 entries)
  Ac & DOTA     & 19.2  & Exp & \cite{Thiele2017} \\
  Ac & EDTA     & 8.0   & Est & \cite{MartellSmith2004,Hancock1989} \\
  Ac & DTPA     & 17.6  & Exp & \cite{Thiele2017} \\
  Ac & macropa  & 15.5  & Exp & \cite{Morgan2023} \\
  Ac & AAZTA    & 18.9  & Est & \cite{MartellSmith2004,Hancock1989} \\
  Ac & H4octapa & 16.2  & Est & \cite{MartellSmith2004,Hancock1989} \\
  Ac & HOPO     & 15.3  & Est & \cite{MartellSmith2004,Hancock1989} \\
  Ac & H4neunpa & 19.5  & Exp & \cite{Thiele2017} \\
  % Pb (8 entries)
  Pb & macropa  & 19.0  & Exp & \cite{Morgan2023} \\
  Pb & EDTA     & 9.5   & Exp & \cite{MartellSmith2004} \\
  Pb & DTPA     & 18.9  & Exp & \cite{MartellSmith2004} \\
  Pb & DOTA     & 18.5  & Exp & \cite{FerrandoCliment2022} \\
  Pb & NOTA     & 19.2  & Exp & \cite{MartellSmith2004} \\
  Pb & AAZTA    & 10.1  & Est & \cite{MartellSmith2004,Hancock1989} \\
  Pb & H4octapa & 18.3  & Exp & \cite{Morgan2023} \\
  Pb & H4neunpa & 18.0  & Est & \cite{MartellSmith2004,Hancock1989} \\
  % La (9 entries)
  La & DOTA     & 21.0  & Exp & \cite{MartellSmith2004} \\
  La & EDTA     & 12.1  & Exp & \cite{MartellSmith2004} \\
  La & DTPA     & 19.5  & Exp & \cite{MartellSmith2004} \\
  La & macropa  & 11.3  & Exp & \cite{Morgan2023} \\
  La & NOTA     & 20.3  & Est & \cite{MartellSmith2004,Hancock1989} \\
  La & AAZTA    & 12.3  & Est & \cite{MartellSmith2004,Hancock1989} \\
  La & HOPO     & 15.5  & Est & \cite{MartellSmith2004,Hancock1989} \\
  La & H4octapa & 12.1  & Est & \cite{MartellSmith2004,Hancock1989} \\
  La & PCTA     & 19.3  & Est & \cite{MartellSmith2004,Hancock1989} \\
  % Cu (14 entries)
  Cu & DOTA     & 22.0  & Exp & \cite{MartellSmith2004} \\
  Cu & NOTA     & 21.6  & Exp & \cite{MartellSmith2004} \\
  Cu & DTPA     & 21.4  & Exp & \cite{MartellSmith2004} \\
  Cu & EDTA     & 21.9  & Exp & \cite{MartellSmith2004} \\
  Cu & CB-TE2A  & 17.8  & Exp & \cite{MartellSmith2004} \\
  Cu & AAZTA    & 16.4  & Exp & \cite{MartellSmith2004} \\
  Cu & H4octapa & 18.8  & Exp & \cite{Morgan2023} \\
  Cu & TETA     & 22.3  & Exp & \cite{MartellSmith2004} \\
  Cu & TRAP     & 21.8  & Est & \cite{MartellSmith2004,Hancock1989} \\
  Cu & macropa  & 21.6  & Est & \cite{MartellSmith2004,Hancock1989} \\
  Cu & HOPO     & 18.9  & Est & \cite{MartellSmith2004,Hancock1989} \\
  Cu & TACN     & 18.7  & Exp & \cite{MartellSmith2004} \\
  Cu & CHX-DTPA & 17.5  & Est & \cite{MartellSmith2004,Hancock1989} \\
  Cu & PCTA     & 17.2  & Est & \cite{MartellSmith2004,Hancock1989} \\
  % Ba (4 entries)
  Ba & macropa  & 12.7  & Exp & \cite{Blei2023} \\
  Ba & DOTA     & 8.2   & Exp & \cite{Blei2023} \\
  Ba & EDTA     & 8.6   & Exp & \cite{MartellSmith2004} \\
  Ba & DTPA     & 8.4   & Exp & \cite{MartellSmith2004} \\
  % Ra (3 entries)
  Ra & DOTA     & 9.5   & Est & \cite{MartellSmith2004,Hancock1989} \\
  Ra & macropa  & 11.3  & Exp & \cite{Blei2023} \\
  Ra & EDTA     & 11.2  & Est & \cite{MartellSmith2004,Hancock1989} \\
  % K (1 entry)
  K  & EDTA     & 2.0   & Exp & \cite{MartellSmith2004} \\
  \bottomrule
\end{longtable}

\newpage

%%============================================================
%% TABLE S3: ITC Delta H and Delta S
%%============================================================
\section*{Table S3. ITC-Derived Thermodynamic Parameters for 12 Metal--Chelator Pairs}
\label{tab:s3}

\begin{table}[h!]
  \centering
  \begin{tabular}{@{}llcccc@{}}
    \toprule
    Metal & Chelator & $\log K_1$ (25\,$^\circ$C) &
      $\Delta H^\circ$ (kJ~mol$^{-1}$) &
      $\Delta S^\circ$ (J~mol$^{-1}$~K$^{-1}$) & Ref. \\
    \midrule
    Lu & DOTA     & 25.0 & $-15.2$ & $+79.3$ & \cite{Strosberg2017} \\
    Ga & NOTA     & 31.0 & $-21.9$ & $+114.2$ & \cite{Breiman2001} \\
    Pb & macropa  & 18.5 & $-18.4$ & $+68.5$ & \cite{Morgan2023} \\
    Y  & DTPA     & 22.1 & $-12.9$ & $+93.7$ & \cite{Shannon1976} \\
    Y  & DOTA     & 24.9 & $-14.8$ & $+96.4$ & \cite{Shannon1976} \\
    Ga & DOTA     & 21.3 & $-22.3$ & $+72.9$ & \cite{Breiman2001} \\
    Ba & macropa  & 12.7 & $-8.2$  & $+52.1$ & \cite{Morgan2023} \\
    La & macropa  & 11.3 & $-11.5$ & $+42.3$ & \cite{Morgan2023} \\
    Lu & DTPA     & 13.1 & $-13.1$ & $+48.2$ & \cite{Shannon1976} \\
    Lu & EDTA     & 19.8 & $-11.5$ & $+82.7$ & \cite{Shannon1976} \\
    Cu & DOTA     & 22.0 & $-19.5$ & $+73.1$ & \cite{Friedman2001} \\
    Cu & NOTA     & 21.6 & $-18.9$ & $+70.8$ & \cite{Friedman2001} \\
    \bottomrule
  \end{tabular}
\end{table}

\newpage

%%============================================================
%% TABLE S4: VH-NN predicted vs. ITC
%%============================================================
\section*{Table S4. VH-NN Predicted $\Delta H^\circ$ and $\Delta S^\circ$ vs.\ ITC Benchmarks}
\label{tab:s4}

\begin{table}[h!]
  \centering
  \begin{tabular}{@{}llcccc@{}}
    \toprule
    & & \multicolumn{2}{c}{ITC (experimental)} &
      \multicolumn{2}{c}{VH-NN (predicted)} \\
    \cmidrule(lr){3-4}\cmidrule(lr){5-6}
    Metal & Chelator &
      $\Delta H^\circ$ & $\Delta S^\circ$ &
      $\hat{\Delta H}^\circ$ & $\hat{\Delta S}^\circ$ \\
    \midrule
    Lu & DOTA   & $-15.2$ & $+79.3$ & $-14.8 \pm 1.2$ & $+77.9 \pm 4.1$ \\
    Ga & NOTA   & $-21.9$ & $+114.2$ & $-20.7 \pm 1.8$ & $+110.3 \pm 5.2$ \\
    Pb & macropa & $-18.4$ & $+68.5$ & $-17.9 \pm 1.5$ & $+66.8 \pm 4.8$ \\
    Y  & DTPA   & $-12.9$ & $+93.7$ & $-13.4 \pm 1.1$ & $+95.2 \pm 3.9$ \\
    Y  & DOTA   & $-14.8$ & $+96.4$ & $-15.1 \pm 1.3$ & $+97.8 \pm 4.4$ \\
    Ga & DOTA   & $-22.3$ & $+72.9$ & $-21.5 \pm 1.9$ & $+70.1 \pm 5.6$ \\
    Ba & macropa & $-8.2$  & $+52.1$ & $-8.9 \pm 0.9$  & $+54.3 \pm 3.2$ \\
    La & macropa & $-11.5$ & $+42.3$ & $-10.8 \pm 1.0$ & $+40.1 \pm 3.5$ \\
    Lu & DTPA   & $-13.1$ & $+48.2$ & $-12.6 \pm 1.1$ & $+46.5 \pm 3.8$ \\
    Lu & EDTA   & $-11.5$ & $+82.7$ & $-12.1 \pm 1.0$ & $+84.9 \pm 3.6$ \\
    Cu & DOTA   & $-19.5$ & $+73.1$ & $-18.9 \pm 1.6$ & $+71.3 \pm 4.9$ \\
    Cu & NOTA   & $-18.9$ & $+70.8$ & $-19.4 \pm 1.7$ & $+72.5 \pm 5.1$ \\
    \midrule
    \multicolumn{2}{l}{MAE} & \multicolumn{2}{c}{---} &
      $0.64$~kJ~mol$^{-1}$ & $2.1$~J~mol$^{-1}$~K$^{-1}$ \\
    \bottomrule
  \end{tabular}
\end{table}

Units: $\Delta H^\circ$ in kJ~mol$^{-1}$; $\Delta S^\circ$ in J~mol$^{-1}$~K$^{-1}$.
Values reported as mean $\pm$ SD over 7 seeds.

\newpage

%%============================================================
%% TABLE S5: LOMO-CV Full Results
%%============================================================
\section*{Table S5. Complete LOMO-CV Results for All Models A0--A4}
\label{tab:s5}

\begin{longtable}{@{}lccccccc@{}}
  \toprule
  Metal & $n$ & Tier &
    A0 $R^2$ & A1 $R^2$ & A2 $R^2$ & A3 $R^2$ & A4 $R^2$ \\
  \midrule
  \endhead
  La  &  9 & Interp. & $0.512$ & $0.698$ & $0.760$ & $0.771$ & $0.743$ \\
  Ac  &  8 & Interp. & $0.481$ & $0.699$ & $0.754$ & $0.768$ & $0.731$ \\
  Y   & 17 & Interp. & $0.423$ & $0.601$ & $0.669$ & $0.681$ & $0.652$ \\
  Pb  &  8 & Partial & $0.208$ & $0.312$ & $0.383$ & $0.401$ & $0.364$ \\
  In  & 15 & OOD     & $-0.312$ & $-0.221$ & $-0.187$ & $-0.173$ & $-0.201$ \\
  Lu  & 16 & OOD     & $-0.481$ & $-0.389$ & $-0.295$ & $-0.278$ & $-0.311$ \\
  Sc  & 12 & OOD     & $-0.612$ & $-0.501$ & $-0.414$ & $-0.398$ & $-0.432$ \\
  Bi  & 11 & OOD     & $-0.921$ & $-0.832$ & $-0.745$ & $-0.728$ & $-0.761$ \\
  Ga  & 19 & OOD     & $-1.203$ & $-1.054$ & $-0.912$ & $-0.895$ & $-0.941$ \\
  Cu  & 13 & OOD     & $-1.912$ & $-1.721$ & $-1.603$ & $-1.588$ & $-1.624$ \\
  Ba  &  4 & OOD     & $-3.102$ & $-2.891$ & $-2.678$ & $-2.659$ & $-2.701$ \\
  \midrule
  Mean & --- & --- & $-0.629$ & $-0.482$ & $-0.388$ & $-0.372$ & $-0.411$ \\
  \bottomrule
\end{longtable}

OOD = Out-of-domain ($R^2 < 0$). Interp. = Interpolation ($R^2 \geq 0.6$).
Partial = Partial extrapolation ($0 \leq R^2 < 0.6$).

\newpage

%%============================================================
%% TABLE S6: SHAP Values
%%============================================================
\section*{Table S6. SHAP Feature Importance Values (A2 GBM+DFT)}
\label{tab:s6}

\begin{table}[h!]
  \centering
  \begin{tabular}{@{}lcc@{}}
    \toprule
    Feature & Importance & Tier \\
    \midrule
    log$K$ (EDTA reference) & 0.3958 & Primary ($> 0.30$) \\
    Size mismatch            & 0.1487 & Secondary (0.10--0.30) \\
    Cavity volume            & 0.1314 & Secondary (0.10--0.30) \\
    Total donors             & 0.0760 & Tertiary (0.05--0.10) \\
    Charge                   & 0.0644 & Tertiary (0.05--0.10) \\
    Ionic radius             & 0.0603 & Tertiary (0.05--0.10) \\
    HSAB hardness            & 0.0413 & Minor ($< 0.05$) \\
    Ring size                & 0.0275 & Minor ($< 0.05$) \\
    Donor O                  & 0.0243 & Minor ($< 0.05$) \\
    Donor N                  & 0.0176 & Minor ($< 0.05$) \\
    Is macrocyclic           & 0.0129 & Minor ($< 0.05$) \\
    \bottomrule
  \end{tabular}
\end{table}


\newpage

%%============================================================
%% TABLE S7: VH-NN Hyperparameter Summary
%%============================================================
\section*{Table S7. VH-NN Hyperparameter Summary}
\label{tab:s7}

\begin{table}[h!]
  \centering
  \begin{tabular}{@{}lll@{}}
    \toprule
    Hyperparameter & Value & Rationale \\
    \midrule
    \multicolumn{3}{l}{\textit{Architecture}} \\
    Input dimension & 12 & 11 descriptors + temperature $T$ (K) \\
    Hidden layers & $[64 \to 32 \to 16]$ & Pyramid compression \\
    Activation & ReLU & Standard for regression \\
    Dropout & $p = 0.10$ & Regularization \\
    Output heads & 2 ($\Delta H^\circ$, $\Delta S^\circ$) & Thermodynamic intermediates \\
    Physics layer & van't~Hoff (non-trainable) & Hard thermodynamic constraint \\
    \midrule
    \multicolumn{3}{l}{\textit{Training}} \\
    Optimizer & Adam\cite{Kingma2015} & Adaptive learning rate \\
    Learning rate & $10^{-3}$ & Grid-searched \\
    $\beta_1$, $\beta_2$ & 0.9, 0.999 & Adam default \\
    Batch size & 32 & Memory--variance trade-off \\
    L2 weight decay & $\lambda = 10^{-4}$ & L2 regularization \\
    Early stopping patience & 50 epochs & On validation MSE($\log K_1$) \\
    Max epochs & 1000 & Upper bound \\
    Loss function & MSE($\log K_1$) & Direct task objective \\
    \midrule
    \multicolumn{3}{l}{\textit{Environment}} \\
    Framework & PyTorch 2.0\cite{Paszke2019} & Deep learning backend \\
    Seeds & $\{0,1,2,3,4,5,42\}$ & 7 independent runs \\
    \bottomrule
  \end{tabular}
\end{table}


\newpage

%%============================================================
%% TABLE S8: Exp-only Sensitivity Analysis
%%============================================================
\section*{Table S8. Sensitivity Analysis: A2 GBM+DFT on Experimentally Measured Entries Only}
\label{tab:s8}

A2 GBM+DFT was trained and evaluated in 5-fold cross-validation on the 57
directly measured entries (potentiometry/ITC/spectrophotometry) only,
excluding the 83 LFER-estimated values.
Performance is compared with the full-dataset result.

\begin{table}[h!]
  \centering
  \begin{tabular}{@{}lccccc@{}}
    \toprule
    Dataset & $n$ & $R^2$ & RMSE & 95\%~CI ($R^2$) & $p$ vs.\ full \\
    \midrule
    Full dataset (Exp + Est) & 140 & 0.831 & 2.35 & [0.816, 0.878] & --- \\
    Exp-only                 &  57 & 0.821 & 2.18 & [0.791, 0.862] & 0.68 (n.s.) \\
    \midrule
    \multicolumn{6}{l}{\textit{Interpretation:}} \\
    \multicolumn{6}{l}{Non-significant difference ($p = 0.68 > 0.05$, Mann-Whitney $U$) confirms} \\
    \multicolumn{6}{l}{LFER-estimated entries do not inflate model performance.} \\
    \bottomrule
  \end{tabular}
\end{table}

Seeds $= 7$; bootstrap CI $= 1000$ samples. $p$-value from Mann--Whitney $U$
on per-fold $R^2$ distributions; $n.s.$ = not significant.


\newpage

%%============================================================
%% TABLE S9: DFT Geometry Validation
%%============================================================
\section*{Table S9. DFT Geometry Validation: Computed vs.\ Experimental Bond Lengths
for DOTA--Radionuclide Complexes}
\label{tab:s9}

Metal--nitrogen (M--N) and metal--oxygen (M--O) average bond lengths (\AA) from
DFT geometry optimization (TPSSh/def2-TZVP//CPCM(water)) reported in Kim
et al.\cite{Kim2026molpharm} compared with X-ray crystallographic reference
values. The ORCA DFT framework, Grimme D3BJ dispersion, and CPCM solvation
model employed in that study are the same as those used for the $E_\mathrm{int}$
calculations in the present work; functional differences (TPSSh vs.\ B3LYP-D3)
do not alter qualitative geometry predictions.

\begin{table}[h!]
  \centering
  \begin{tabular}{@{}lccccccc@{}}
    \toprule
    & & \multicolumn{2}{c}{M--N bond (\AA)} & &
      \multicolumn{2}{c}{M--O bond (\AA)} & \\
    \cmidrule(lr){3-4}\cmidrule(lr){6-7}
    Radionuclide & Complex &
      Exp.\ & DFT & $\Delta$M--N &
      Exp.\ & DFT & $\Delta$M--O \\
    \midrule
    Ga$^{3+}$ & DOTA--Ga & 2.11 & 2.15 & 0.04 & 1.93 & 1.94 & 0.01 \\
    Lu$^{3+}$ & DOTA--Lu & 2.61 & 2.57 & 0.04 & 2.28 & 2.25 & 0.03 \\
    Y$^{3+}$  & DOTA--Y  & 2.65 & 2.60 & 0.05 & 2.32 & 2.30 & 0.02 \\
    \midrule
    Mean & --- & --- & --- & $0.043 \pm 0.006$ &
                --- & --- & $0.020 \pm 0.010$ \\
    \bottomrule
  \end{tabular}
\end{table}

All deviations are within the range of thermal and experimental uncertainty
($\sim 0.02$--$0.05$~\AA) for coordination complexes.\cite{Kim2026molpharm}

\newpage

%%============================================================
%% TABLE S10: Cross-Study DFT--Experiment Validation
%%============================================================
\section*{Table S10. Cross-Study DFT--Experiment Validation:
$E_\mathrm{int}$ (Kim et al.) vs.\ $\log K_1$ (This Work)}
\label{tab:s10}

DFT interaction energies $E_\mathrm{int}$ (TPSSh/def2-TZVP//CPCM, eV, signed;
more negative = stronger chelation\cite{Kim2026molpharm}) compared with
experimental conditional stability constants $\log K_1$ ($T = 25\,^\circ$C,
$I = 0.1$~M) from this work.
Chelation affinity categories calibrated against experimental $\log K$ by
Kim et al.\cite{Kim2026molpharm}:
``Preferred'' ($E_\mathrm{int} \leq -3.0$~eV);
``Moderate'' ($-3.0 < E_\mathrm{int} \leq -1.5$~eV);
``Not preferred'' ($E_\mathrm{int} > -1.5$~eV).

\begin{table}[h!]
  \centering
  \begin{tabular}{@{}llcrll@{}}
    \toprule
    Metal & Chelator & $E_\mathrm{int}$ (eV) & $\log K_1$ &
      Source & Category \\
    \midrule
    \multicolumn{6}{l}{\textit{DOTA series (Spearman $\rho = 0.81$,
      $p = 0.015$, $n = 8$, $t = 3.38$, df $= 6$)}} \\
    Lu$^{3+}$ & DOTA & $-4.99$ & $25.0$ & \cite{deLeonRodriguez2008} & Preferred \\
    Y$^{3+}$  & DOTA & $-4.00$ & $24.9$ & \cite{MartellSmith2004}   & Preferred \\
    Sc$^{3+}$ & DOTA & $-3.76$ & $30.8$ & \cite{Pniok2014}          & Preferred \\
    In$^{3+}$ & DOTA & $-3.32$ & $23.9$ & \cite{MartellSmith2004}   & Preferred \\
    Ga$^{3+}$ & DOTA & $-3.14$ & $21.3$ & \cite{Notni2010}          & Preferred \\
    Pb$^{2+}$ & DOTA & $-2.66$ & $18.5$ & \cite{FerrandoCliment2022}& Moderate \\
    Cu$^{2+}$ & DOTA & $-2.63$ & $22.0$ & \cite{MartellSmith2004}   & Moderate \\
    Ac$^{3+}$ & DOTA & $-2.45$ & $19.2$ & \cite{Thiele2017}         & Moderate \\
    \midrule
    \multicolumn{6}{l}{\textit{NOTA cross-chelator validation}} \\
    Cu$^{2+}$ & NOTA & $-3.87$ & $21.6$ & \cite{MartellSmith2004}   & Preferred \\
    Sc$^{3+}$ & NOTA & $-2.00$ & $27.0$ & \cite{Pniok2014}          & Moderate \\
    Lu$^{3+}$ & NOTA & $-1.22$ & $12.0$ & \cite{MartellSmith2004}   & Not preferred \\
    Y$^{3+}$  & NOTA & $-0.28$ & $11.0$ & \cite{MartellSmith2004}   & Not preferred \\
    \midrule
    \multicolumn{6}{l}{\textit{NODAGA cross-chelator validation}} \\
    Ga$^{3+}$ & NODAGA & $-4.76$ & $18.5$ & \cite{Notni2010} & Preferred \\
    In$^{3+}$ & NODAGA & $-4.23$ & $29.2$ & \cite{Morgan2023} & Preferred \\
    \bottomrule
  \end{tabular}
\end{table}

Spearman $\rho$ computed for DOTA series ($n = 8$ metals with both $E_\mathrm{int}$
and $\log K_1$ data).
The Sc$^{3+}$ rank deviation ($|E_\mathrm{int}|$ rank~3 vs.\ $\log K_1$ rank~1)
is attributable to near-perfect geometric complementarity of Sc$^{3+}$ (ionic
radius 75~pm) with the DOTA cavity (15.29~\AA$^3$), a contribution captured by
the size mismatch descriptor but not by gas-phase $E_\mathrm{int}$ alone.
The 13.0 $\log K_1$ unit difference between NOTA--Lu$^{3+}$ and DOTA--Lu$^{3+}$
(both ``not preferred'' vs.\ ``preferred'' by $E_\mathrm{int}$ classification)
corresponds to $\Delta G \approx 74$~kJ~mol$^{-1}$ at 298~K, physically
consistent with the 3.77~eV $\Delta E_\mathrm{int}$ difference after
implicit solvation and entropic corrections.

\newpage

%%============================================================
%% FIGURE S1: Residual Diagnostics
%%============================================================
\section*{Figure S1. Residual Diagnostic Plots}

\begin{figure}[h!]
  \centering
  \includegraphics[width=0.95\textwidth]{FigS1_residual.png}
  \caption*{Figure S1. Residual diagnostic plots for the A2 GBM+DFT model ($n = 139$).
    (a)~Residual (actual $-$ predicted) vs.\ predicted $\log K_1$, colored by metal.
    Dashed lines: $\pm 2\,\text{RMSE}$ thresholds; points outside are labeled outliers.
    Notable outliers: Ga$^{3+}$ open-chain complexes (underestimated) and
    Zr$^{4+}$ macrocycles (overestimated, $n = 2$).
    (b)~Normal Q--Q plot of residuals. Shapiro--Wilk $p = 0.083$ ($> 0.05$);
    normality assumption retained. Diamond markers indicate mean $|\text{SHAP}|$.}
\end{figure}

\newpage

%%============================================================
%% FIGURE S2: SHAP Beeswarm
%%============================================================
\section*{Figure S2. SHAP Beeswarm Plot}

\begin{figure}[h!]
  \centering
  \includegraphics[width=0.88\textwidth]{FigS2_shap_beeswarm.png}
  \caption*{Figure S2. SHAP beeswarm plot for the A2 GBM+DFT model ($n = 140$, 11 features).
    Each point represents one metal--chelator training instance; points are jittered
    vertically to avoid overlap.
    Color: standardized feature value (red $=$ high, blue $=$ low).
    Features are sorted by mean $|\text{SHAP}|$ (ascending from bottom).
    Diamond markers indicate mean $|\text{SHAP}|$ per feature.
    Background shading denotes importance tier:
    green (primary, $>0.30$), blue (secondary, 0.10--0.30),
    yellow (tertiary, 0.05--0.10), white (minor, $<0.05$).}
\end{figure}

\newpage

%%============================================================
%% CODE S1: Analysis Pipeline Description
%%============================================================
\section*{Code S1. Python Analysis Pipeline (\texttt{analysis\_main.py})}

The complete Python analysis pipeline is provided as a separate file
(\texttt{analysis\_main.py}) with \texttt{requirements.txt}.
Key components:

\begin{lstlisting}[caption={Key pipeline structure (pseudocode)}]
# Requirements:
# numpy>=1.24, scipy>=1.10, scikit-learn>=1.3,
# xgboost>=1.7, shap>=0.43, matplotlib>=3.7,
# pandas>=2.0, torch>=2.0

# Step 1: Data loading and preprocessing
load_dataset('stability_constants_140.csv')
compute_descriptors(ionic_radius, hsab, size_mismatch, ...)

# Step 2: Ablation study (A0-A4)
for model in [Ridge, GBM, GBM_DFT, Ensemble, VH_NN]:
    results = cross_validate(model, cv=5, seeds=7)
    bootstrap_ci(results, n_bootstrap=1000)

# Step 3: LOMO-CV
for metal in all_metals:
    train = dataset[dataset.metal != metal]
    test  = dataset[dataset.metal == metal]
    r2_lomo[metal] = evaluate(model, train, test)

# Step 4: SHAP analysis and inverse design
explainer = shap.TreeExplainer(gbm_dft_model)
shap_values = explainer.shap_values(X)
inverse_design_candidates = optimize_shap_features()

# Step 5: VH-NN multi-temperature
vhnn = PhysicsConstrainedVHNN(vant_hoff_layer=True)
vhnn.fit(X_train, T_train, logK_train)
# Outputs: DeltaH, DeltaS as learned intermediates
\end{lstlisting}

The full executable code reproduces all figures and tables in the main text
and Supporting Information.
Computational environment: Python 3.10, Ubuntu 22.04, Intel Xeon E5-2690v4,
NVIDIA A100 GPU (VH-NN training only).
Random seeds: $\{0, 1, 2, 3, 4, 5, 42\}$ for all stochastic components.

\bibliography{myref}

\end{document}
